\subsection*{The Constitutive Role of the Soul}

Aristotle's account of the soul in \textit{De Anima} reveals its constitutive role in the actualization chain with particular clarity. The soul is ``a substance in the sense of the form of a natural body'' (Aristotle, \textit{On The Soul (De Anima)}, p.~18). As form, the soul is the first actuality of a natural body that has life potentially; it is what makes the body into a living organism capable of perceiving, moving, and thinking. The soul does not exist apart from the body---it is the body's organizing principle, its way of being what it is. Hence the soul is itself an instance of the priority of actuality over potentiality: it is the form that must be actual before the body's capacities can be exercised. This point connects directly to the motion--time parallel. Just as motion is the actuality of the potential qua potential---the process by which what can be comes to be---so the soul is the actuality of the body's potential for life (Aristotle, \textit{On The Soul (De Anima)}, p.~18). And just as time is the numbered aspect of motion, so the soul is the numbering principle that articulates temporal flow. The parallel is not merely analogical; it reflects the deep structural isomorphism between the physical and psychical domains of Aristotle's ontology. Motion and time are the physical ground; soul and perception are the psychical development. $A_0$ and $A_1$ are two stages of a single process of actualization.

Aristotle's hierarchical ordering of psychic capacities in \textit{De Anima} II.3 reinforces this point. Among mortal beings, ``those which possess calculation have all the other powers above mentioned, while the converse does not hold---indeed some live by imagination alone, while others have not even imagination'' (Aristotle, \textit{On The Soul (De Anima)}, p.~22). The hierarchy of capacities---nutritive, perceptive, imaginative, rational---mirrors the chain of actualities. Each level presupposes the actualization of the levels below it. Perception presupposes the nutritive capacity (for only living things perceive); imagination presupposes perception (for \textit{phantasmata} are the residues of perceptual motions); thought presupposes imagination (for the soul never thinks without a phantasm). Similarly, each psychic capacity presupposes the physical actualities of motion and time at $A_0$: without motion there is no alteration of the sense-organ, and without time there is no sequential unfolding of perceptual, imaginative, and cognitive acts.

Here the architectural weight of $A_0$ becomes most explicit. Aristotle asks at \textit{Physics} IV.14, 223a25--27 whether time could exist if soul did not, and his answer is carefully qualified: ``But if nothing but soul, or in soul reason, is qualified to count, it is impossible for there to be time unless there is soul, but only that of which time is an attribute---i.e., if movement can exist without soul.'' The conditional is precise. Motion, Aristotle allows, may exist without soul; natural bodies move whether or not anyone observes them. But time---which is the \textit{numbered} aspect of motion---seems to require a counter, and where there is no counter there can be no number counted. Thus while motion does not depend on soul, time does: time is real as the countable structure of motion but becomes fully actual as time only in relation to a soul capable of counting.

This chapter adopts a stronger interpretive reading of Aristotle's conditional. On the reading advanced here, time without a counting soul is not merely unnumbered but \textit{ontologically incomplete}---present as countability, structurally real as the before-and-after of motion, yet not fully actualized as time. The interpretive move goes beyond what Aristotle explicitly states. Aristotle's own formulation is epistemically cautious (if-there-is-no-counter-then-no-number); the stronger ontological claim---that time's full actuality requires soul---is the reading this chapter takes up. The reading is defensible because it coheres with Aristotle's broader priority-of-actuality commitments: if time is fully actual only as numbered, then what makes the numbering possible---the soul---is constitutive of time's full actuality, not merely its epistemic availability.

This interpretive reading supports the architectural claim that $A_0$ is not a neutral substrate but an actuality whose full actualization is achieved only with the arrival of the soul at $A_1$. The chain, on this reading, is not a sequence of stages running sequentially upon an already-complete foundation; it is an iterated structure in which each stage is what makes the prior stage what it fully is. $A_0$ is structurally present from the beginning---motion and time are co-eternal, always already in act---but time's full actualization as numbered requires the soul that $A_1$ introduces. The soul, in turn, exists only because motion and time have already made natural beings possible. The interdependence is reciprocal; the stages are not linearly ordered but nested and co-constitutive.

The decisive architectural payoff lies in the parallel between motion's mode of being and time's mode of being. Motion, as \textit{energeia ateles}, occupies an intermediate ontological register: neither pure potentiality nor complete actuality, but the active exercise of the potential \textit{as} potential. Time, on the interpretive reading just advanced, occupies a structurally analogous register: neither brute succession nor fully articulated numerical order, but the countable aspect of motion that is on its way to being counted. Motion is actual as ongoing unfolding; time is actual as countable structure. Neither is a static, self-contained \textit{energeia}; both are ongoing processes whose very incompleteness is what allows them to ground further actualities.

The parallel licenses everything downstream. Because motion is \textit{energeia ateles}, it can serve as the unmoved originator of the motions that produce completed actualities at $A_1$, $A_2$, and beyond---it is ongoing enough to be the source, incomplete enough not to terminate in itself. Because time is the countable structure of motion, it can become fully articulated only as the soul actualizes the counting---it is real enough to condition sequential ordering, open enough to be developed as the numbered continuum in which the subsequent motions unfold. $A_0$'s mode of being is thus not static grounding but \textit{active grounding}: the ongoing process whose very incompleteness makes further completion possible.

Every $A_n$ in the chain inherits this intermediate-register character. Each completed actuality at $A_1$, $A_2$, $A_3$ is both the terminus of its antecedent motion and the unmoved originator of its successor---an \textit{entelecheia} in relation to what precedes it, an \textit{arche} in relation to what follows. The chain is iterated because each of its nodes has the same dual structure that $A_0$ establishes as primitive: completion looking backward, potentiality looking forward, actuality all the way down. The fundamental motion of the chain---what Aristotle might call its \gk{καθόλου} structure---is not sequential passage from stage to stage but the repeated exercise of a single ontological pattern: \textit{energeia ateles} generating \textit{entelecheia} generating the conditions for a new \textit{energeia ateles}. $A_0$ is where this pattern first emerges and from which the chain inherits its primitive structure.

The mandatory theoretical synthesis with Heidegger's temporal analysis in \textit{Being and Time} becomes relevant at precisely this juncture. The \gk{φάντασμα} as memory-image---which functions simultaneously as an object of contemplation in itself and as a likeness (\gk{εἰκών}) of the thing remembered---presupposes the temporal continuum established at $A_0$. Memory requires the before and after that time provides; the phantasma's dual aspect as present image and past likeness is intelligible only within a temporal manifold. Aristotle's treatment of time in \textit{Physics} IV.11 directly prepares the ground for this, and Heidegger's analysis of temporality in \textit{Being and Time} can be understood, at least in part, as a radicalization of the Aristotelian insight that time is constitutively bound up with the soul's activity (Multiple Authors, \textit{Heidegger and Rhetoric}, pp.~124--125). The innovations in Heidegger's theory of time represent ``the very fundamental commitment in Heidegger's theory to a temporal manifold, rather than to a linear sequence of past, present, and future `nows'\,'' (Multiple Authors, \textit{Heidegger and Rhetoric}, pp.~124--125). Thus the Aristotelian $A_0$---motion and time as first actuality---is not merely a historical curiosity but a living philosophical resource whose implications Heidegger himself recognized and radicalized. Chalmers's work on the virtual and the real offers a further perspective on the constitutive role of the perceiving subject. The lucid dreamer constitutes a virtual world through representational activity (Chalmers, \textit{Reality+}, pp.~38--40), and this observation, while situated within a contemporary framework, echoes the Aristotelian insight that the soul's constitutive activity shapes the character of the perceptual world; the virtual world of the lucid dreamer is structured by the dreamer's representational capacities, just as the actual world of the perceiver is structured by the soul's capacity to number motion and receive sensible forms.
